% !TeX root = ../papers/nc-p3.tex
\section{一文字次数の移動による全通常階数15窓の九次自由性}
\label{sec:rank-fifteen-ninth-character}
\begin{theorem}\label{thm:rank-fifteen-ninth-character}
標数零の正の固定曲線の全通常階数15窓で、四始点に共通の四つの中心四次式を選び、
九次まで中心単項式基底を取れる。各始点の九次経路は220次元、中心商は20次元である。
\end{theorem}
\begin{proof}
第\ref{sec:rank-fifteen-ninth-orbit-obstruction}節で固定四列の九次が23次元となった
障害を、別の四列の二変数消去なしに処理する。残った標数3の単元域で、正規化前の
36成分のWalsh表示を用いる。そのKlein文字次数を $0,1,2,3$ とする。
全四位置で文字に同じ $a$ を加える単項行列による自己同型が $a=1,2$ で存在する。
実際、第\ref{sec:rank-fifteen-genus-two-unit-line}節のWalsh和で
四文字を $(i,j,k,t)$、$h=i\oplus k=j\oplus t$ とする。
$1$ の加算では $i\oplus j$ と $k,t$ の上位ビットが保たれ、各項が同じである。
$2$ の加算では和の添字 $u$ と $\beta u$ の上位ビットの和が符号に加わる。
$\beta=(0,3,2,1)$ なのでこれは $u$ の下位ビット、すなわち $h$ の下位ビットである。
位置0,2の各文字にその下位ビットの符号を掛ければ元に戻る。
四周期対角基底変更による正規化後も、文字の置換は同じである。
長さ4では合計文字は変わらないので文字0の四次中心空間 $V$ は保たれ、
奇数長では合計文字を任意に平行移動できる。

まず固定した $c^2=X$ の六表示を調べる。固定四列 $(1,2,4,5)$ の
長さ8までの始点0の計算に続き、長さ9では文字3の関係行列だけを消去する。
この文字の商次元は5である。残る三文字の上界はここでは用いない。
\path{computations/verify_rank_fifteen_ninth_character.py} はこの計算で逆数を取った
主元を集める。既知のチャートの単元を除いた $G(v,c)$ の二次数は $(181,160)$ である。
$B=Y/X^3$、$(v^3-v)^2=B$ の六基底における係数の最大公約元は1である。
第\ref{sec:rank-fifteen-oriented-eighth-cut}節の符号を保つ終結式上界の次数は67である。
$c$ の次数2以下は既知で、それ以外に残る $(c,B)$ のFrobenius軌道は
三次が八つ、四次が三つである。各軌道で元の窓に戻して十三次まで全四始点の
標準商を検算する。従ってその十一軌道では既出の全次数の持上げを使える。

その他では、固定した $c$ の少なくとも一表示で、固定四平面 $U\subset V$ に対し
$\dim (A/(U))_{0,9,3}\le5$ である。まず\emph{同じ剰余窓の中で}四平面を動かす。
上の自己同型によりどの文字 $h$ でも $\dim(A/(U_h))_{0,9,h}\le5$ の四平面が存在する。
各不等式は $\operatorname{Gr}(4,V)$ 上で開く。中心が文字0なので選び直した全ての
四平面でも商に文字次数があり、各文字の関係挿入行列の小行列式が使える。
四つの非空開集合を交差させれば、同一の四平面で九次の全商が20次元以下となる。
この操作は固定四列のまま九次が20になることを必要としない。

五つの中心の基底が持ち上がるので、得た剰余四平面の係数も持ち上げられる。
$c$ を指定した六表示のため、$\bar p=c^3$ により目的の添字保持窓で始点0の上界を得る。
他の始点にも移動した窓で同じ議論を行い、標数零の全四次中心のGrassmann多様体で
四始点の開条件と八次までの自由性の開条件を交差させる。
中心倍による生成と未切断経路の一般点からの下界から
\[
  220=\binom{12}{3}\ \le\ \dim A_{p,9}\
     \le\ \dim Q_{p,9}+4\cdot40+10\cdot4\ \le\ 220 .
\]
よって全て等号であり、九次でも中心商基底の持上げと中心単項式の積が基底になる。
\end{proof}
